\documentclass[11pt]{article}

\usepackage[letterpaper,margin=1in]{geometry}
\usepackage{fontspec}
\setmainfont{DejaVu Serif}
\setsansfont{DejaVu Sans}
\setmonofont{DejaVu Sans Mono}[Scale=0.86]
\usepackage{microtype}
\usepackage{amsmath,amsthm,amssymb,mathtools}
\usepackage{booktabs,array,tabularx,longtable}
\usepackage{enumitem}
\usepackage{xcolor}
\usepackage{hyperref}
\usepackage{cleveref}
\usepackage{fvextra}
\usepackage{tcolorbox}
\usepackage{caption}

\definecolor{dkblue}{HTML}{173B57}
\definecolor{dkred}{HTML}{8B1E2D}
\definecolor{dkgray}{HTML}{F3F5F7}
\definecolor{dkline}{HTML}{CBD3DA}
\hypersetup{
  colorlinks=true,
  linkcolor=dkblue,
  citecolor=dkblue,
  urlcolor=dkblue,
  pdftitle={A counterexample to Davis--Kahan Proposition 4.4 with related work}
}

\newtheorem{theorem}{Theorem}[section]
\newtheorem{proposition}[theorem]{Proposition}
\newtheorem{lemma}[theorem]{Lemma}
\newtheorem{remark}[theorem]{Remark}
\newtheorem{question}[theorem]{Question}
\theoremstyle{definition}
\newtheorem{definition}[theorem]{Definition}

\newcommand{\Hh}{\mathcal H}
\newcommand{\R}{\mathbb R}
\newcommand{\Id}{I}
\newcommand{\trnorm}[1]{\lVert #1\rVert_1}
\newcommand{\opnorm}[1]{\lVert #1\rVert_{\mathrm{op}}}
\newcommand{\spanof}{\operatorname{span}}
\newcommand{\svals}{\operatorname{sv}}
\newcommand{\DK}{Davis--Kahan}
\newcommand{\lean}{\textsc{Lean}}

\DefineVerbatimEnvironment{LeanCode}{Verbatim}{
  fontsize=\footnotesize,
  breaklines=true,
  breakanywhere=true,
  frame=single,
  rulecolor=\color{dkline},
  framesep=3mm,
  xleftmargin=2mm,
  xrightmargin=2mm,
  baselinestretch=1.05
}

\newtcolorbox{claimbox}{
  colback=dkred!4,
  colframe=dkred!75!black,
  boxrule=0.8pt,
  arc=1.5mm,
  left=3mm,right=3mm,top=2.5mm,bottom=2.5mm
}
\newtcolorbox{statusbox}{
  colback=dkblue!4,
  colframe=dkblue!70!black,
  boxrule=0.7pt,
  arc=1.5mm,
  left=3mm,right=3mm,top=2.5mm,bottom=2.5mm
}

\title{\textbf{A Counterexample to Davis--Kahan Proposition 4.4 (Draft technical note for expert review)}\\[0.45em]
\large A finite-dimensional trace-norm obstruction, its Lean certification,\\
and the semantic alignment required to interpret the formal result}

\author{Jonathan Crall \and Edward Yang \and GPT-5.6 \and GPT-5.5 \and Fable 5 \and Opus 4.8}

\date{22 July 2026}

\begin{document}
\maketitle

\begin{abstract}
We claim that Proposition 4.4 of Chandler Davis and W. M. Kahan's 1970 paper
\emph{The Rotation of Eigenvectors by a Perturbation. III} is false as printed.
The proposition asserts, in a real Hilbert space and under the angle bound
$\Theta\leq \pi/3$, that the direct rotation minimizes
$\lVert I-V\rVert$ among all orthogonal operators carrying one subspace onto
another, simultaneously for every unitarily invariant norm.  We give an
explicit counterexample in $\R^4$ for the trace norm.  The two principal angles
are both $\pi/4$, but an admissible orthogonal competitor has trace displacement
$2\sqrt2$, whereas the direct rotation has trace displacement
$4\sqrt{2-\sqrt2}>2\sqrt2$.  The example has been certified in Lean 4.

A machine-checked counterexample to a formal statement does not by itself prove
that a historical prose statement has been refuted: a human must establish the
semantic correspondence between the source notation and the formal objects.
Accordingly, this note records both the mathematical counterexample and an
explicit semantic-alignment audit.  It also identifies equation (4.3) in the
published proof as a failing step: for the present example, its left-hand side is
$2\sqrt2$ while its asserted lower bound is $4$.

The broader norm-dependent phenomenon was partly known: Bolotnikov, Li, and
Rodman showed in 2001 that canonical minimization can fail for Schatten
$p$-norms with $1\leq p<2$.  Their explicit obstruction, however, occurs at the
right-angle endpoint and does not refute the Davis--Kahan short-angle claim.
Moreover, Qiu, Zhang, and Li repeated the $\pi/3$ assertion in 2006.  We are
unaware of an earlier published counterexample within the advertised
$\theta_{\max}\leq\pi/3$ regime; the historical priority claim is therefore
stated cautiously.
\end{abstract}

\begin{claimbox}
\textbf{Primary claim.}
\emph{Davis--Kahan Proposition 4.4 is false as stated in the original published
paper.}
\end{claimbox}

\tableofcontents

\section{Statement of the issue}

The paper \cite{DavisKahan1970} studies the geometry of two closed subspaces
$P\Hh$ and $Q\Hh$ and the canonical \emph{direct rotation} $U$ carrying the
first onto the second.  Its Section 4 concerns extremal properties of $U$.
The original article explicitly notes that Section 4 is not required for the
rest of the paper; thus the issue discussed here does not undermine the later
single-angle, double-angle, or numerical perturbation theorems.

On printed page 22, Proposition 4.4 states, in the notation of the paper, that
if $V$ is a unitary taking $P\Hh$ onto $Q\Hh$ in a real Hilbert space and if
$\Theta\leq \pi/3$, then $\lVert I-V\rVert$ is minimized, for every
unitarily invariant norm, when $V=U$, the direct rotation.  In modern finite-
dimensional language, the claim is the following.

\begin{proposition}[Finite-dimensional specialization of \DK{} Proposition 4.4]
\label{prop:published-claim}
Let $E$ be a finite-dimensional real inner-product space, let $\mathcal U$ and
$\mathcal V$ be subspaces whose largest principal angle is at most $\pi/3$, and
let $R_{\mathcal U,\mathcal V}$ denote the direct rotation from $\mathcal U$ to
$\mathcal V$.  For every orthogonal operator $W$ satisfying
$W\mathcal U=\mathcal V$ and every unitarily invariant norm $N$ on
$\operatorname{End}(E)$,
\[
 N(I-R_{\mathcal U,\mathcal V})\leq N(I-W).
\]
\end{proposition}

A counterexample to this finite-dimensional specialization is automatically a
counterexample to the published Hilbert-space proposition.

\section{Related work and historical positioning}
\label{sec:related}

The present result should be distinguished from the previously known fact that
canonical or direct rotations need not solve every global unitary-completion
problem for every unitarily invariant norm.  The new point asserted here is
narrower: the published Davis--Kahan conclusion fails \emph{inside} its stated
real short-angle regime.  Our present historical assessment is:

\begin{statusbox}
General failures of direct-rotation optimality for Schatten $p<2$ norms were
known by 2001, but the published examples lie outside the Davis--Kahan
short-angle regime.  To our knowledge, Proposition~4.4 has not previously been
observed to fail under its stated hypothesis
$\theta_{\max}\leq\pi/3$.
\end{statusbox}

This is deliberately not an unconditional priority claim.  Bibliographic
absence is difficult to prove, particularly for older, non-digitized material,
private correspondence, and subscriber-only citation databases.  The evidence
nevertheless supports the more precise statement that no widely recognized or
indexed correction of the short-angle claim has been found.

\subsection{The original claim and its advertised repair}

The 1970 paper does not merely assert a bound for the operator norm, the square
norm, or the restriction of the displacement to the source subspace.  It claims
that, in a real Hilbert space and under $\Theta\leq\pi/3$, the direct rotation
minimizes $\lVert I-V\rVert$ for \emph{every} unitarily invariant norm
\cite{DavisKahan1970}.  The trace norm is therefore explicitly included.  The
surrounding discussion contrasts Proposition~4.4 with two preceding examples:
one showing failure for a sufficiently large real angle, and another showing
failure in a complex space.  The reality assumption and the $\pi/3$ bound are
presented as repairing those failures.  There is no hypothesis of a single
principal plane, distinct principal angles, or a competitor preserving the
principal-plane decomposition.

\subsection{Historical treatments of the CS decomposition and direct rotation}

Paige and Wei's historical account of the cosine--sine decomposition reviews
the contributions of Davis and Kahan and explains the relation between the CS
decomposition and the direct rotations of Davis and Kato
\cite{PaigeWei1994}.  It is a natural place where a known error in the geometric
extremal theory might have been mentioned, but it does not flag
Proposition~4.4.  This is not proof that the error was unknown in 1994; it is
supporting evidence that no correction had become part of the standard history.

\subsection{The norm-class divide in Bolotnikov--Li--Rodman}

Bolotnikov, Li, and Rodman studied the minimization problem
\[
 \min\{\lVert Z-I\rVert_\Phi:
       Z\text{ unitary and }Z\mathcal M_j=\mathcal N_j\}
\]
for chains of subspaces and unitarily invariant norms
\cite{BolotnikovLiRodman2001}.  Their formulation includes, as the one-subspace
case, the ambient completion problem considered in Davis--Kahan Section~4.
They introduce the class of \emph{$Q$-norms} and prove canonical minimization
for that class.  For Schatten norms, the $Q$-norm condition holds exactly for
\[
  \lVert\cdot\rVert_p,\qquad 2\leq p\leq\infty,
\]
because $\lVert A\rVert_p^2=\lVert A^*A\rVert_{p/2}$.  This encompasses the
Frobenius norm and the operator norm and parallels the passage from the valid
squared-displacement theorem of Davis--Kahan Proposition~4.3.

The same paper explicitly shows that the canonical result is generally false
for Schatten norms $1\leq p<2$; see its Lemma~3.3.  Thus failure for the trace
norm was known in substance by 2001.  However, the explicit negative example is
obtained by setting the cosine parameter in its CS block to $a=0$.
Geometrically,
\[
 a=\cos\theta=0,
 \qquad \theta=\frac\pi2,
\]
so the example is the orthogonal-subspace endpoint and lies completely outside
the Davis--Kahan range $\theta\leq\pi/3$.  The authors conclude that the
canonical theorem is generally false for non-$Q$-norms, but do not investigate
whether a short-angle hypothesis restores it.

This distinction is essential.  The present note does not claim to discover
that trace-like norms can favor a noncanonical completion.  It supplies an
explicit example showing that the particular $\pi/3$ repair proposed by Davis
and Kahan also fails.

\subsection{The 2006 near miss}

Qiu, Zhang, and Li revisited unitarily invariant geometry on Grassmann spaces
\cite{QiuZhangLi2006}.  In their discussion of extremal properties of the
direct rotation, they repeat three conclusions attributed to Davis--Kahan:
minimal restricted displacement, minimal squared displacement, and minimal
full displacement when the largest canonical angle is at most $\pi/3$.  The
third is precisely the claim refuted here.  They describe the $\pi/3$
restriction as undesirable, but do not question its correctness, and instead
prove the different, globally valid statement that a direct rotation minimizes
a logarithmic or geodesic displacement $\lVert\log Z\rVert_\Phi$.

The paper contains an especially close near miss.  In a majorization example it
compares the angle vectors
\[
 \left(\frac\pi4,\frac\pi4\right)
 \quad\text{and}\quad
 \left(\frac\pi2,0\right)
\]
and observes the reversal
\[
 \sin\frac\pi4+\sin\frac\pi4
   =\sqrt2
   >1
   =\sin\frac\pi2+\sin0.
\]
This is essentially the same ``spread versus concentrate'' mechanism used in
our counterexample: two equal moderate rotations can cost more in an additive
trace-like norm than one larger rotation together with one fixed plane.  The
2006 authors apply the observation to a Grassmann metric and do not connect it
to the ambient completion problem $\min\lVert I-Z\rVert_1$.

The historical evidence is stronger because Chi-Kwong Li coauthored both the
2001 paper establishing general failures for non-$Q$-norms and the 2006 paper
repeating the $\pi/3$ assertion.  The 2006 paper also acknowledges helpful
discussions with Chandler Davis.  The natural interpretation is that the
general norm-class obstruction was understood, while the short-angle claim was
still believed to be correct.  This is evidence rather than proof of what the
authors knew.

\subsection{Later structural and norm-dependent literature}

Dou, Shi, Cui, and Du give general explicit descriptions of intertwining
operators and direct rotations of two orthogonal projections
\cite{DouShiCuiDu2017}.  Their work improves structural existence and
representation results associated with Kato, Avron--Seiler--Simon, and
Davis--Kahan, but it does not report a correction to the Section~4 extremal
statement.

Norm-dependent optimality also appears in orthogonal Procrustes problems.  Cape
studies how an optimizer depends on the chosen norm
\cite{Cape2020}.  Such problems optimize frame alignments of the form
$\inf_Z\lVert X-YZ\rVert$ and reinforce the general lesson that a canonical
optimizer may depend on the norm.  They are nevertheless different from the
ambient orthogonal-extension problem
\[
 \inf\{\lVert I-Z\rVert:Z\mathcal U=\mathcal V\}
\]
and do not identify a Davis--Kahan error.

Books, surveys, and later applications found in this review tend to rely on the
safer neighboring statements: minimization of $(I-Z)P$, minimization of the
squared displacement $(I-Z)^*(I-Z)$ in arbitrary unitarily invariant norms, or
minimality for the operator and Frobenius norms.  These claims remain compatible
with both the $Q$-norm theory and the present counterexample.

\subsection{Errata search and the official publication record}

SIAM's current author guidance states that an accepted author erratum is
published as a separate item and attached to the original article's PDF; a
production correction is likewise appended to the original PDF
\cite{SIAMErrataPolicy}.  We found no erratum, correction, expression of
concern, or attached notice on the official article record for DOI
\texttt{10.1137/0707001}, and targeted searches did not locate a SIAM-recorded
correction of Proposition~4.4.  This supports the narrow statement:

\begin{quote}
There is no SIAM-recorded correction of Proposition~4.4 that we could locate.
\end{quote}

It does not exclude an unpublished observation, correspondence, or a correction
buried in unrelated literature.

\subsection{Search limitations and archival route}

The literature search included the exact $\pi/3$ assertion, ``Proposition 4.4''
with Davis--Kahan, direct rotation combined with trace, nuclear, and Schatten
norms, the angle pattern $(\pi/4,\pi/4)$ versus $(\pi/2,0)$, the exact algebraic
costs in the example, and terms such as \emph{erratum}, \emph{correction},
\emph{false}, and \emph{counterexample}.  The only directly relevant hits were
the general $p<2$ failure in 2001 and the 2006 repetition of the short-angle
claim.

The main unresolved bibliographic limitations are subscriber-only forward
citation lists, MathSciNet and zbMATH reviews that were not exhaustively
exported, dissertations and non-digitized notes, and private correspondence.
The University of Toronto holds the Chandler Davis fonds, comprising 59 boxes
of research material and extensive correspondence, with most files open to
researchers \cite{ChandlerDavisFonds}.  The most promising archival targets are
correspondence and drafts from approximately 1967--1971, exchanges with Kahan,
notes for \emph{The Rotation of Eigenvectors by a Perturbation. III}, and later
correspondence concerning unitary-invariant norms and direct rotations.  An
archival request is a genuinely independent historical route that could alter
the priority assessment.

\subsection{Chronology and confidence assessment}

\begin{center}
\begin{tabularx}{0.97\textwidth}{>{\raggedright\arraybackslash}p{0.17\textwidth}
                                  >{\raggedright\arraybackslash}p{0.29\textwidth}
                                  >{\raggedright\arraybackslash}X}
\toprule
\textbf{Source} & \textbf{What it establishes} & \textbf{Relation to this counterexample}\\
\midrule
Davis--Kahan (1970) & All-UI-norm minimization in real spaces for $\Theta\leq\pi/3$ & The false statement itself.\\
Paige--Wei (1994) & History of the CS decomposition and direct rotation & Does not flag Proposition~4.4.\\
Bolotnikov--Li--Rodman (2001) & Positive theory for $Q$-norms; failure for Schatten $p<2$ & Explicit negative example is at $\theta=\pi/2$, not in the short-angle regime.\\
Qiu--Zhang--Li (2006) & Repeats the $\pi/3$ assertion; proves logarithmic metric optimality & Strong evidence that the precise error was still unnoticed; contains a closely related majorization example.\\
Dou--Shi--Cui--Du (2017) & General formulas for direct rotations and intertwining operators & Structural theory, with no correction of the extremal statement.\\
Cape (2020) & Norm-dependent behavior in a Procrustes problem & Different feasible-set geometry; no Davis--Kahan correction.\\
Present note & Real $\R^4$ counterexample with angles $(\pi/4,\pi/4)$ & Refutes Proposition~4.4 inside its advertised regime.\\
\bottomrule
\end{tabularx}
\end{center}

Our confidence levels are as follows:
\begin{itemize}
\item \textbf{High confidence:} there is no widely recognized or mainstream
  published correction of Proposition~4.4.
\item \textbf{Moderately high confidence:} no indexed paper gives an explicit
  counterexample under $\theta_{\max}\leq\pi/3$.
\item \textbf{Moderate confidence:} the real $\R^4$, $(\pi/4,\pi/4)$,
  trace-norm counterexample is mathematically new.
\item \textbf{Low confidence:} no mathematician ever noticed the error
  privately, in obscure literature, or in correspondence.
\end{itemize}

The strongest responsible novelty language is therefore:

\begin{statusbox}
General failures of direct-rotation optimality for Schatten $p<2$ norms were
previously known, but the published examples lie outside the Davis--Kahan
short-angle regime.  To our knowledge, Proposition~4.4 has not previously been
observed to fail under its stated hypothesis
$\theta_{\max}\leq\pi/3$.
\end{statusbox}

Before submission, the highest-value remaining due diligence is an
institutional MathSciNet/zbMATH/Scopus citation export, an archival request to
the Chandler Davis fonds, and direct outreach to experts or surviving authors.
Further ordinary web searching is unlikely to change the assessment as much as
those three routes.


\section{The explicit counterexample}
\label{sec:counterexample}

We now give the exact configuration used in the Lean development.  Let
$E=\R^4$ with standard orthonormal basis $e_0,e_1,e_2,e_3$, and set
\[
 \mathcal U=\spanof\{e_0,e_1\}.
\]
Define the orthogonal matrix
\begin{equation}
\label{eq:W}
 W=\frac12
 \begin{pmatrix}
  1&-1&-1&-1\\
  1& 1& 1&-1\\
 -1&-1& 1&-1\\
  1&-1& 1& 1
 \end{pmatrix},
 \qquad
 \mathcal V=W\mathcal U.
\end{equation}
The four rows of the sign matrix are mutually orthogonal and have squared norm
$4$, so $W^TW=I$.  Thus $W$ is an admissible orthogonal competitor carrying
$\mathcal U$ onto $\mathcal V$ by construction.

\subsection{The principal angles are both $\pi/4$}

Let $E_{\mathcal U}$ be the $4\times2$ matrix with columns $e_0,e_1$ and let
$E_{\mathcal V}=WE_{\mathcal U}$.  The cross-Gram matrix is
\[
 E_{\mathcal U}^T E_{\mathcal V}
 =E_{\mathcal U}^TWE_{\mathcal U}
 =\frac12
 \begin{pmatrix}
  1&-1\\
  1& 1
 \end{pmatrix}
 =\frac1{\sqrt2}
 \begin{pmatrix}
  \cos(\pi/4)&-\sin(\pi/4)\\
  \sin(\pi/4)& \cos(\pi/4)
 \end{pmatrix}.
\]
Its two singular values are both $1/\sqrt2$.  Hence the two principal angles
between $\mathcal U$ and $\mathcal V$ are
\[
 \theta_1=\theta_2=\arccos(1/\sqrt2)=\frac\pi4.
\]
In particular,
\[
 \theta_{\max}=\frac\pi4<\frac\pi3,
\]
so the angle hypothesis of Proposition 4.4 is satisfied strictly.  The pair is
also acute: neither principal angle is $\pi/2$.

\subsection{Identification of the direct rotation}

The target projection is
\[
 P_{\mathcal V}
 =\frac12
 \begin{pmatrix}
 1&0&0&1\\
 0&1&-1&0\\
 0&-1&1&0\\
 1&0&0&1
 \end{pmatrix},
 \qquad
 P_{\mathcal U}=\operatorname{diag}(1,1,0,0).
\]
The canonical intertwiner used to define the direct rotation is
\[
 S=P_{\mathcal V}P_{\mathcal U}
   +(I-P_{\mathcal V})(I-P_{\mathcal U}).
\]
A direct calculation gives
\[
 S^TS=\frac12 I.
\]
Therefore the unitary polar factor of $S$ is $R=\sqrt2S$, namely
\begin{equation}
\label{eq:R}
 R=\frac1{\sqrt2}
 \begin{pmatrix}
  1& 0& 0&-1\\
  0& 1& 1& 0\\
  0&-1& 1& 0\\
  1& 0& 0& 1
 \end{pmatrix}.
\end{equation}
This is the direct rotation from $\mathcal U$ to $\mathcal V$.  It rotates the
principal plane $\spanof\{e_0,e_3\}$ through $\pi/4$ and the principal plane
$\spanof\{e_1,e_2\}$ through $\pi/4$ (up to orientation).

For later use, one obtains the particularly simple displacement identity
\begin{equation}
\label{eq:Rgram}
 (I-R)^T(I-R)=(2-\sqrt2)I.
\end{equation}
Thus all four singular values of $I-R$ are
$\sqrt{2-\sqrt2}$, and therefore
\begin{equation}
\label{eq:Rtrace}
 \trnorm{I-R}=4\sqrt{2-\sqrt2}.
\end{equation}

\subsection{Trace displacement of the competitor}

Introduce the orthonormal basis
\[
 \begin{aligned}
 m_0&=\frac{e_0+e_2}{\sqrt2},&
 m_1&=\frac{e_1+e_3}{\sqrt2},\\
 m_2&=\frac{e_0-e_2}{\sqrt2},&
 m_3&=\frac{e_1-e_3}{\sqrt2}.
 \end{aligned}
\]
The matrix $W$ acts by
\[
 Wm_0=m_1,\qquad Wm_1=-m_0,\qquad Wm_2=m_2,\qquad Wm_3=m_3.
\]
Hence, in this basis,
\[
 W\simeq Q_{\pi/2}\oplus I_2,
\]
where $Q_{\pi/2}$ is planar rotation through $\pi/2$.  The singular values of
$I-W$ are consequently
\[
 \sqrt2,\ \sqrt2,\ 0,\ 0,
\]
and
\begin{equation}
\label{eq:Wtrace}
 \trnorm{I-W}=2\sqrt2.
\end{equation}

\subsection{The strict contradiction}

Combining \cref{eq:Rtrace,eq:Wtrace},
\[
 \trnorm{I-W}=2\sqrt2
 <4\sqrt{2-\sqrt2}
 =\trnorm{I-R}.
\]
For completeness, both sides are positive, and squaring reduces the strict
inequality to
\[
 8<16(2-\sqrt2),
\]
or equivalently $\sqrt2<3/2$, which follows from $2<9/4$.

We have therefore proved the following.

\begin{theorem}[Counterexample to \DK{} Proposition 4.4]
\label{thm:counterexample}
There exist acute two-dimensional subspaces $\mathcal U,\mathcal V\subset\R^4$
with both principal angles equal to $\pi/4$, and an orthogonal operator $W$
with $W\mathcal U=\mathcal V$, such that
\[
 \trnorm{I-W}<\trnorm{I-R_{\mathcal U,\mathcal V}}.
\]
Consequently Proposition 4.4 is false for the trace norm, despite the strict
angle bound $\theta_{\max}<\pi/3$.
\end{theorem}

\section{Multiplicity concentration at a repeated angle}

The example is not an isolated numerical accident.  It is a manifestation of
multiplicity mixing.  A convenient coordinate-free family is obtained in
$\R^2\oplus\R^2$ by setting
\[
 \mathcal U_\theta
 =\left\{\frac1{\sqrt2}(x,x):x\in\R^2\right\},
 \qquad
 W_\theta=Q_{2\theta}\oplus I,
 \qquad
 \mathcal V_\theta=W_\theta\mathcal U_\theta.
\]
The cross-Gram operator is
\[
 \frac{I+Q_{2\theta}}2=\cos\theta\,Q_\theta,
\]
so both principal angles are $\theta$.  The direct rotation uses two principal
planes and distributes the movement as $(\theta,\theta)$, whereas the competitor
concentrates it as $(2\theta,0)$.

For a planar rotation $Q_\alpha$, both singular values of $I-Q_\alpha$ equal
$2\sin(\alpha/2)$.  Consequently
\[
 \trnorm{I-R_\theta}=8\sin(\theta/2),
 \qquad
 \trnorm{I-W_\theta}=4\sin\theta
 =8\sin(\theta/2)\cos(\theta/2).
\]
For every $0<\theta<\pi$, the competitor is strictly cheaper in trace norm.
Thus no positive upper-angle threshold can rescue the unrestricted trace-norm
claim when repeated principal angles permit mixing of the multiplicity space.
The choice $\theta=\pi/4$ yields the exact matrix model in
\cref{sec:counterexample}, up to an orthogonal change of coordinates.

The same construction also explains why the operator-norm statement behaves
differently.  The direct rotation has largest displacement singular value
$2\sin(\theta/2)$, while the concentrated competitor has largest singular value
$2\sin\theta>2\sin(\theta/2)$.  Concentration decreases an additive Schatten
norm such as the trace norm but increases the maximum singular value.

\section{The apparent failure in the published proof}
\label{sec:proofgap}

The proof of Proposition 4.4 invokes equation (4.3).  For an even integer $\nu$
and mutually orthogonal two-dimensional principal-plane projections
$\Omega_k$, that equation asserts, in the paper's notation,
\begin{equation}
\label{eq:paper43}
 \lVert K\rVert_\nu
 \geq
 \sum_{k=1}^{\nu/2}\lVert K\Omega_k\rVert_2.
\end{equation}
Here $\lVert\cdot\rVert_j$ denotes the Ky Fan $j$-norm, not the Hilbert--Schmidt
norm.  Orthogonality of the domains $\Omega_kE$ does not imply orthogonality of
the ranges $K\Omega_kE$, and without such output orthogonality the trace norm of
$K\sum_k\Omega_k$ need not decompose into a sum of trace norms.

The present counterexample directly falsifies \cref{eq:paper43}.  Put
$K=I-W$, $\nu=4$, and take the direct-rotation principal planes
\[
 \Omega_1E=\spanof\{e_0,e_3\},
 \qquad
 \Omega_2E=\spanof\{e_1,e_2\}.
\]
A direct calculation gives
\[
 (K\Omega_1)^T(K\Omega_1)=I_2,
 \qquad
 (K\Omega_2)^T(K\Omega_2)=I_2.
\]
Thus each $K\Omega_k$ has singular values $(1,1)$ and
\[
 \lVert K\Omega_1\rVert_2+\lVert K\Omega_2\rVert_2=2+2=4.
\]
On the other hand,
\[
 \lVert K\rVert_4=\trnorm{I-W}=2\sqrt2<4.
\]
Therefore equation (4.3), in the generality in which it is used, is false.
This pinpoints a concrete logical failure in the printed proof rather than merely
exhibiting a theorem whose proof resists formalization.

\section{Lean formulation}
\label{sec:lean}

The formal development used for this note is the repository state at commit
\texttt{816e91a}.  The counterexample is in the file
\begin{center}
\small\path{DavisKahan/FiniteDimensional/DirectRotation/ShortRotationCounterexample.lean}
\end{center}
The full commit hash is recorded in the appendix.  The code below is reproduced
in a form suitable for semantic review.

\subsection{A Lean statement corresponding to the published positive claim}

The following is a finite-dimensional formal statement corresponding to
Proposition 4.4.  It is displayed as a proposition-valued definition rather than
as a theorem, because the assertion is false.

\begin{LeanCode}
def DavisKahanProposition4_4_Finite : Prop :=
  ∀ (E : Type*)
    [NormedAddCommGroup E]
    [InnerProductSpace ℝ E]
    [FiniteDimensional ℝ E]
    (U V : Submodule ℝ E)
    [U.HasOrthogonalProjection]
    [V.HasOrthogonalProjection]
    (hacute : IsAcute U V)
    (hshort : principalAngles U V 0 ≤ Real.pi / 3)
    (W : E ≃ₗᵢ[ℝ] E)
    (hmap : U.map W.toLinearMap = V)
    (N : UnitarilyInvariantNorm ℝ E),
      N (LinearMap.id -
          (directRotation U V hacute).toLinearMap) ≤
        N (LinearMap.id - W.toLinearMap)
\end{LeanCode}

This is the natural finite-dimensional specialization needed for refutation.
The explicit `hacute` argument is not an additional mathematical restriction:
the source hypothesis $\Theta\leq\pi/3$ already excludes a principal angle
$\pi/2$, and the Lean constructor for the canonical direct rotation requests a
proof of acuteness.

\subsection{The machine-checked counterexample theorem}

The repository contains the following theorem.

\begin{LeanCode}
theorem shortRotation_fullDisplacement_refuted :
    ∃ (U V : Submodule ℝ (EuclideanSpace ℝ (Fin 4)))
      (hacute : IsAcute U V)
      (W : EuclideanSpace ℝ (Fin 4) ≃ₗᵢ[ℝ]
        EuclideanSpace ℝ (Fin 4)),
      U.map W.toLinearMap = V ∧
      principalAngles U V 0 ≤ Real.pi / 3 ∧
      kyFanSum 4 (LinearMap.id - W.toLinearMap) <
        kyFanSum 4
          (LinearMap.id -
            (directRotation U V hacute).toLinearMap) :=
  ⟨U4, V4, acute, Wequiv, rfl,
    principalAngle_le, kyFanSum_lt⟩
\end{LeanCode}

The key exact-value lemmas are:

\begin{LeanCode}
theorem kyFanSum_displacement_W :
    kyFanSum 4 (LinearMap.id - Wlin) =
      2 * Real.sqrt 2

 theorem kyFanSum_displacement_R :
    kyFanSum 4
      (LinearMap.id -
        (directRotation U4 V4 acute).toLinearMap) =
      4 * Real.sqrt (2 - Real.sqrt 2)
\end{LeanCode}

The file also certifies orthogonality of $W$, acuteness of the pair, the principal
angle bound, the singular values of both displacements, and the final strict
inequality.

\section{Semantic alignment between prose and Lean}
\label{sec:alignment}

A formal theorem is only as relevant to a historical claim as the interpretation
connecting the two.  Lean certifies implications among formal definitions; it
does not certify that those definitions are the intended meanings of prose,
notation, or a scanned source.  Establishing that correspondence is a distinct
mathematical and historiographic obligation.

\begin{longtable}{>{\raggedright\arraybackslash}p{0.22\textwidth}
                  >{\raggedright\arraybackslash}p{0.31\textwidth}
                  >{\raggedright\arraybackslash}p{0.37\textwidth}}
\toprule
\textbf{Published concept} & \textbf{Lean representation} & \textbf{Alignment argument}\\
\midrule
\endfirsthead
\toprule
\textbf{Published concept} & \textbf{Lean representation} & \textbf{Alignment argument}\\
\midrule
\endhead
Real Hilbert space
& `EuclideanSpace ℝ (Fin 4)`
& $\R^4$ is a real Hilbert space.  A finite-dimensional counterexample refutes a theorem stated for arbitrary real Hilbert spaces.\\
\addlinespace
Subspaces $P\Hh$ and $Q\Hh$
& `U V : Submodule ℝ E`
& In finite dimension, closed subspaces and linear submodules coincide for the purposes at issue; orthogonal projections exist automatically.\\
\addlinespace
Unitary in a real space
& `E ≃ₗᵢ[ℝ] E`
& A surjective real linear isometry is exactly an orthogonal operator.  The matrix $W$ is proved to preserve inner products and have an inverse.\\
\addlinespace
$V(P\Hh)=Q\Hh$
& `U.map W.toLinearMap = V`
& `Submodule.map` is the image of the source subspace under the underlying linear map.  Equality is stronger than mere inclusion.\\
\addlinespace
The direct rotation $U$ of Davis--Kahan
& `directRotation U V hacute`
& The Lean object is the unitary polar factor of the canonical intertwiner $P_VP_U+P_{V^\perp}P_{U^\perp}$.  In the example this polar factor is the explicit matrix $R$ in \cref{eq:R}; it has the expected principal-plane action and carries $\mathcal U$ onto $\mathcal V$.\\
\addlinespace
Angle condition $\Theta\leq\pi/3$
& `principalAngles U V 0 ≤ Real.pi / 3`
& `principalAngles` is formed from the singular values of the directed sine map and sorted in decreasing order.  Index $0$ is the largest angle.  Independently, the cross-Gram calculation shows both angles equal $\pi/4$, so the conclusion does not depend on a delicate indexing convention.\\
\addlinespace
Every unitarily invariant norm
& `N : UnitarilyInvariantNorm ℝ E`
& The formal structure permits seminorm-level objects, so it is slightly broader than the paper's class of norms.  This does not affect the refutation: `kyFanSum 4` in dimension four is the genuine trace/nuclear norm, a standard unitarily invariant norm.\\
\addlinespace
Full displacement $I-V$
& `LinearMap.id - W.toLinearMap`
& This is the full ambient displacement.  It is not the restricted displacement $(I-W)P_U$ and not the positive square $(I-W)^*(I-W)$.  Those nearby statements have different validity.\\
\bottomrule
\end{longtable}

\subsection{Why human semantic review is necessary}

The central alignment risks are the following.

\begin{enumerate}[leftmargin=2.2em]
\item \textbf{Source fidelity.}  One must inspect the original printed paper,
not merely a modern transcription, to verify the quantifiers, the real-space
restriction, the angle threshold, the competitor class, and the use of the full
displacement.  The original page contains no additional hypothesis excluding
multiplicity mixing.

\item \textbf{Notation overload.}  The paper uses $U$ for the direct rotation
and $V$ for a general competitor.  The formalization uses submodule variables
`U`, `V` and a competitor `W`.  Variable names are semantically irrelevant, but
a human must avoid mistaking the symbols.

\item \textbf{Angle conventions.}  Davis--Kahan's angle operator on the full
ambient decomposition duplicates nonzero principal angles, whereas the finite
Lean list is one-sided.  The maximum angle is nevertheless the same.  In this
example the independent cross-Gram computation removes any possible doubt.

\item \textbf{Norm conventions.}  The paper's subscripted norms
$\lVert\cdot\rVert_\nu$ are Ky Fan norms, while its square norm uses separate
notation.  The Lean `UnitarilyInvariantNorm` structure is formulated at the
seminorm level, but the witness is the genuine full Ky Fan $4$-sum, exactly the
trace norm in $\R^4$.  Thus the slightly broader formal class cannot be the
source of the counterexample.

\item \textbf{Canonical-rotation conventions.}  Several equivalent definitions
of direct rotation occur in the literature.  The formal object must be checked
against the paper's positivity/principal-square-root characterization, not merely
called a `directRotation`.  The polar-factor computation and explicit matrix
$R$ provide this check in the example.
\end{enumerate}

Thus the formal theorem establishes a rigorous counterexample to the Lean
statement.  The conclusion that it refutes the published Proposition 4.4 rests on
the preceding human-verified interpretation.  That interpretation is unusually
strong here because every ingredient can be represented by explicit $4\times4$
matrices and checked independently of the formal library.

\section{What remains valid}

The counterexample distinguishes three superficially similar extremal claims.

\begin{enumerate}[leftmargin=2.2em]
\item The direct rotation minimizes the \emph{restricted} displacement
$(I-W)P_{\mathcal U}$ in the relevant singular-value and unitarily invariant
norm senses.  The competitor cannot exploit arbitrary behavior on
$\mathcal U^\perp$ in that problem.

\item The positive displacement square
$(I-W)^*(I-W)$ satisfies a different majorization/minimality theorem.  Applying
a norm after squaring is not equivalent to applying the trace norm before
squaring.

\item The operator-norm full-displacement minimality is compatible with the
counterexample.  The competitor has a larger largest displacement singular
value, even though the sum of its singular values is smaller.
\end{enumerate}

The false claim is therefore specifically the simultaneous full-displacement
minimality for \emph{every} unitarily invariant norm.

\section{Status and suggested independent checks}

\begin{statusbox}
The counterexample and its supporting algebra are machine-checked.  The
literature reviewed in \cref{sec:related} establishes that the general
$p<2$ failure was already known, while no earlier short-angle refutation was
found.  We therefore use ``apparently new'' and ``to our knowledge'' rather
than asserting unconditional priority.  Before publication, the remaining due
diligence should include MathSciNet and zbMATH citation-chain searches,
inspection of reviews and dissertations, and direct consultation with experts
in matrix nearness and Grassmann geometry.
\end{statusbox}

An expert reviewer can verify the mathematical claim with the following short
checklist.

\begin{enumerate}[leftmargin=2.2em]
\item Confirm $W^TW=I$ for \cref{eq:W}.
\item Compute $E_{\mathcal U}^TWE_{\mathcal U}$ and verify its two singular
values are $1/\sqrt2$.
\item Compute $P_{\mathcal V}$, the canonical intertwiner $S$, and verify
$S^TS=\frac12I$ and $R=\sqrt2S$.
\item Verify \cref{eq:Rgram}, hence
$\trnorm{I-R}=4\sqrt{2-\sqrt2}$.
\item Change to the $m$-basis and verify
$W\simeq Q_{\pi/2}\oplus I_2$, hence
$\trnorm{I-W}=2\sqrt2$.
\item Compare the two exact values.
\item Inspect printed pages 20--23 of \cite{DavisKahan1970}, especially
Proposition 4.4 and equation (4.3), to confirm source alignment.
\end{enumerate}

\appendix

\section{Compact matrix certificate}

For ease of independent reproduction, the entire counterexample can be reduced
to the following identities:
\[
 U=\spanof\{e_0,e_1\},\qquad V=WU,
\]
with $W$ as in \cref{eq:W}, and
\[
 W^TW=I,
 \qquad
 \svals(E_U^TWE_U)=\left(\frac1{\sqrt2},\frac1{\sqrt2}\right),
\]
\[
 R=\frac1{\sqrt2}
 \begin{pmatrix}
  1& 0& 0&-1\\
  0& 1& 1& 0\\
  0&-1& 1& 0\\
  1& 0& 0& 1
 \end{pmatrix},
 \qquad
 (I-R)^T(I-R)=(2-\sqrt2)I,
\]
\[
 \svals(I-W)=(\sqrt2,\sqrt2,0,0),
 \qquad
 \svals(I-R)=
 (\sqrt{2-\sqrt2},\sqrt{2-\sqrt2},
  \sqrt{2-\sqrt2},\sqrt{2-\sqrt2}).
\]
These identities imply the contradiction immediately.

\section{Repository provenance}

The formal result was inspected in the source archive corresponding to the
following commit:
\begin{LeanCode}
816e91a1193c0b3ece24d4ad68ae35059297e521
\end{LeanCode}
The principal declaration is
\[
 \texttt{ForMathlib.DavisKahanTheory.
 shortRotation\_fullDisplacement\_refuted}.
\]
The proof file contains no admitted proof in the counterexample chain and
records exact singular-value and Ky Fan-sum computations for both the competitor
and the direct rotation.

\begin{thebibliography}{99}

\bibitem{BolotnikovLiRodman2001}
V. Bolotnikov, C.-K. Li, and L. Rodman,
\newblock Minimal distortion problems for classes of unitary matrices,
\newblock \emph{Electronic Journal of Linear Algebra}, 8:26--46, 2001.
\newblock \href{https://doi.org/10.13001/1081-3810.1058}{doi:10.13001/1081-3810.1058}.

\bibitem{Cape2020}
J. Cape,
\newblock Orthogonal Procrustes and norm-dependent optimality,
\newblock \emph{Electronic Journal of Linear Algebra}, 36:158--168, 2020.
\newblock \href{https://doi.org/10.13001/ela.2020.5009}{doi:10.13001/ela.2020.5009}.

\bibitem{DavisKahan1970}
C. Davis and W. M. Kahan,
\newblock The rotation of eigenvectors by a perturbation. III,
\newblock \emph{SIAM Journal on Numerical Analysis}, 7(1):1--46, 1970.
\newblock \href{https://doi.org/10.1137/0707001}{doi:10.1137/0707001}.


\bibitem{PaigeWei1994}
C. C. Paige and M. Wei,
\newblock History and generality of the CS decomposition,
\newblock \emph{Linear Algebra and its Applications}, 208--209:303--326, 1994.
\newblock \href{https://doi.org/10.1016/0024-3795(94)90446-4}{doi:10.1016/0024-3795(94)90446-4}.

\bibitem{DouShiCuiDu2017}
Y.-N. Dou, W.-J. Shi, M.-M. Cui, and H.-K. Du,
\newblock General explicit descriptions for intertwining operators and direct rotations of two orthogonal projections,
\newblock \emph{Linear Algebra and its Applications}, 531:575--591, 2017.
\newblock \href{https://doi.org/10.1016/j.laa.2017.06.036}{doi:10.1016/j.laa.2017.06.036}.

\bibitem{SIAMErrataPolicy}
Society for Industrial and Applied Mathematics,
\newblock Information for journal authors: Errors in published articles and publication of errata,
\newblock accessed 22 July 2026.
\newblock \url{https://epubs.siam.org/journal-authors}.

\bibitem{ChandlerDavisFonds}
University of Toronto Archives and Records Management Services,
\newblock Chandler Davis fonds, UTA 1218, 1941--2022, 59 boxes,
\newblock finding aid accessed 22 July 2026.
\newblock \url{https://discoverarchives.library.utoronto.ca/index.php/chandler-davis-fonds}.

\bibitem{QiuZhangLi2006}
L. Qiu, Y. Zhang, and C.-K. Li,
\newblock Unitarily invariant metrics on the Grassmann space,
\newblock \emph{SIAM Journal on Matrix Analysis and Applications},
27(2):507--531, 2005; published online 2006.
\newblock \href{https://doi.org/10.1137/040607605}{doi:10.1137/040607605}.

\end{thebibliography}

\end{document}
